% Generated draft: requires mathematical and source review.
\documentclass[11pt]{article}
\usepackage[margin=25mm]{geometry}
\usepackage{fontspec}
\setmainfont{Latin Modern Roman}
\usepackage{amsmath,amssymb,mathtools,unicode-math}
\setmathfont{Latin Modern Math}
\usepackage{enumitem,needspace,xurl,hyperref}
\setlength{\parindent}{0pt}
\setlength{\parskip}{5pt}
\setlength{\emergencystretch}{3em}
\newcommand{\N}{\mathbb{N}}
\newcommand{\Z}{\mathbb{Z}}
\newcommand{\Q}{\mathbb{Q}}
\newcommand{\R}{\mathbb{R}}
\newcommand{\C}{\mathbb{C}}
\newcommand{\F}{\mathbb{F}}
\newcommand{\A}{\mathbb{A}}
\newcommand{\PP}{\mathbb{P}}
\newcommand{\E}{\mathbb{E}}
\newcommand{\eps}{\varepsilon}
\newcommand{\dd}{\,\mathrm{d}}
\begin{document}
\section*{UnsolvedMath ID: 30000019}
\textbf{Source problem code:} catalogue entry
\subsection{Definitions and mathematical statement}
Let $\Gamma$ be a connected locally finite graph, meaning every vertex has finitely many neighbors, and let $G\le\operatorname{Aut}(\Gamma)$ act transitively on vertices. A $G$-invariant block system $\mathcal B$ is a partition of the vertex set such that each $g\in G$ permutes its blocks. It has finite blocks when every $B\in\mathcal B$ is finite. Put
\[
 K=\{g\in G:gB=B\text{ for every }B\in\mathcal B\},\qquad
 \bar G=G/K\le\operatorname{Sym}(\mathcal B).
\]
The natural quotient-action interpretation of ``finite induced stabilizer'' is
\[
 |\bar G_B|<\infty\quad(B\in\mathcal B),
\]
where $\bar G_B$ is a block stabilizer in the induced permutation action. By transitivity, finiteness for one block is equivalent to finiteness for all.

The conjectural alternative in the verified source statement is: for every such pair $(\Gamma,G)$, either there exists a finite-block imprimitivity system with finite induced stabilizer, or $\Gamma$ is hyperbolic. It places no explicit growth assumption on $\Gamma$ and does not require $G$ to be the full automorphism group.

The excerpt does not define ``hyperbolic'' in the intended graph-theoretic context, nor specify whether singleton block systems are allowed. If Gromov hyperbolicity were intended, it would mean that the unit-edge path metric has $\delta$-thin geodesic triangles for some $\delta\ge0$. That is an explicit possible interpretation, not a verified identification of the report's terminology. The precise hyperbolicity notion and quotient-stabilizer convention must be recovered from the supplied Trofimov reference or report before this becomes a fully formal dichotomy; substituting a stronger familiar notion could change the conjecture.

The literature assessment checked on 21 August 2026 reports strong structural theorems in polynomial-growth cases and says the full hyperbolic-or-finite-block alternative is unresolved. This is dated curation, not an independently established current status or permission to restrict the stated graph class to polynomial growth.
\subsection{Short English statement}
The conjecture asks whether every locally finite vertex-transitive graph has either a finite-block quotient action with finite stabilizers or the report's hyperbolicity property. The exact meaning of hyperbolicity and the allowed block systems must be checked in the source.
\subsection{Sources}
\begin{itemize}
\item \emph{Mini-Workshop: Amalgams for Graphs and Geometries} (2005), Oberwolfach Reports, DOI 10.4171/owr/2004/25, pp. 1311--1342, \url{https://doi.org/10.4171/owr/2004/25}; proposed year supplied as 2004.
\item Literature assessment checked 21 August 2026: Trofimov, \url{https://doi.org/10.1016/j.ejc.2005.05.010}; Tessera--Tointon, \url{https://arxiv.org/abs/1908.06044}.
\item Local UnsolvedMath registry, permanent ID $30000019$, exact source code \texttt{catalogue entry}, created 25 August 2026 and updated 16 September 2026. The statement is labelled corrected and verified with symbols and title hyphenation normalized, and the dichotomy retained.
\end{itemize}
\end{document}
